\documentclass{umk-matematika}

\begin{document}

\maketitlepage
    {Практическая работа №6}
    {Корень n-й степени. Свойства}
    {}

\section{Фундамент (1--4 балла)}

\textbf{Задача 1 (1 балл).} Вычислите: $\sqrt{49}$

\textbf{Задача 2 (2 балла).} Вычислите: $\sqrt[3]{27}$

\textbf{Задача 3 (3 балла).} Вычислите: $\sqrt{36 \cdot 25}$

\textbf{Задача 4 (4 балла).} Вычислите: $\sqrt{\dfrac{81}{9}}$

\section{Стены (5--7 баллов)}

\textbf{Задача 5 (5 баллов).} Вынесите множитель из-под знака корня: $\sqrt{50}$

\textbf{Задача 6 (6 баллов).} Упростите: $\sqrt{27} - \sqrt{12}$

\textbf{Задача 7 (7 баллов).} Вычислите: $\sqrt[3]{-8} + \sqrt[3]{64}$

\section{Кровля (8--10 баллов)}

\textbf{Задача 8 (8 баллов).} Упростите выражение: $\sqrt{x^4}$, если $x < 0$.

\textbf{Задача 9 (9 баллов).} Упростите: $(\sqrt{5} + \sqrt{3})(\sqrt{5} - \sqrt{3})$

\textbf{Задача 10 (10 баллов).} Площадь квадратного участка земли равна 144 м². Найдите длину стороны участка.

\newpage
\section*{Ответы}

\begin{enumerate}
  \item 7.
  \item 3.
  \item 30.
  \item 3.
  \item $5\sqrt{2}$
  \item $\sqrt{3}$
  \item 2.
  \item $x^2$
  \item 2.
  \item сторона участка = 12 м.
\end{enumerate}

\end{document}